\section{Clamp the exact sum before storing}\label{pcm-saturation}
For an exact sum $s$, the output is
\[
 \operatorname{clip}_{16}(s)=
 \begin{cases}
 -32768 & s<-32768,\\
 s & -32768\leq s\leq32767,\\
 32767 & s>32767.
 \end{cases}
\]
The scalar implementation compares against the upper and lower bounds while the
exact sum is still available. On a clipping path it replaces the value and increments
the clip counter once. On an in-range path it stores the exact sum without incrementing.
Both paths then consume one sample.

\begin{table}[htbp]
\centering
\caption{Endpoint and overflow examples. Clipped counts strict out-of-range sums, not the number of output samples equal to an endpoint.}
\tagpdfsetup{table/header-rows={1}}
\begin{tabular}{rrrrr}
\toprule
Input A & Input B & Exact sum & Output & Clipped\\
\midrule
30000 & 2767 & 32767 & 32767 & 0\\
32767 & 1 & 32768 & 32767 & 1\\
30000 & 10000 & 40000 & 32767 & 1\\
-30000 & -2768 & -32768 & -32768 & 0\\
-32768 & -1 & -32769 & -32768 & 1\\
-30000 & -10000 & -40000 & -32768 & 1\\
32767 & -32768 & -1 & -1 & 0\\
\bottomrule
\end{tabular}
\end{table}

A stored endpoint is ambiguous: it can come from a valid exact sum or a clipped
larger magnitude. The return count therefore needs evidence about the arithmetic,
not a later scan counting endpoint values. It also counts samples rather than
clipping directions; an overflow in one sample contributes one, whether positive
or negative.
